\section{The tightening operator}\label{sec:foundations}

Consider
\begin{equation}\label{eq:problem}
 f^*=\min\{f(x):x\in X\cap B_0\},
\end{equation}
where $B_0=[\ell^0,u^0]\subset\R^n$ is a bounded box and $X$ includes
the original constraints and any integrality restrictions. A relaxation
built on $B\subseteq B_0$ has a projected objective
$\phi_B:B\longrightarrow\R\cup\{+\infty\}$: for a lifted formulation,
$\phi_B(x)$ is the infimum of the relaxed objective over lifted points
above $x$, with value $+\infty$ if none exist. Assume
\begin{align}
 \phi_B(x)&\le f(x) &&(x\in X\cap B),\label{eq:validity}\\
 \phi_{B'}(x)&\ge\phi_B(x) &&(B'\subseteq B,\ x\in B').\label{eq:monotonicity}
\end{align}
The second condition is a hypothesis on the whole relaxation construction,
including any rows retained or discarded. It is not a consequence of
calling a relaxation ``McCormick'' if unrelated strengthening rows change
between calls. Assume that the nonempty projected sublevel sets used
below are compact. This holds for the bounded lifted linear relaxations
in the exact reference implementation.

For a finite objective cutoff $U$, define
\begin{equation}\label{eq:operator}
 K_U(B)=\{x\in B:\phi_B(x)\le U\},\qquad
 T_U(B)=\hull K_U(B),
\end{equation}
where $\hull$ is the smallest axis-aligned closed box containing its
argument. The empty set maps to the empty set. Exact Jacobi OBBT computes
the coordinate minima and maxima on one frozen relaxation, then rebuilds
the relaxation on the resulting box. A sequential variant may rebuild
after each coordinate update; this difference matters for rate equalities.

\begin{lemma}[Containment and order]\label{lem:order}
Every feasible point in $X\cap B$ with objective at most $U$ belongs to
$T_U(B)\subseteq B$. If $B'\subseteq B$, then
$T_U(B')\subseteq T_U(B)$. If $U'\le U$, then
$T_{U'}(B)\subseteq T_U(B)$.
\end{lemma}
\begin{proof}
Validity puts the stated feasible points in $K_U(B)$. The box hull stays
in $B$. For $x\in K_U(B')$, monotonicity gives
$\phi_B(x)\le\phi_{B'}(x)\le U$, so $K_U(B')\subseteq K_U(B)$;
taking box hulls preserves inclusion. The cutoff assertion follows
directly from sublevel-set inclusion.
\end{proof}

If a feasible incumbent supplies $U\ge f^*$, exact iterates
$B_{k+1}=T_U(B_k)$ are nested and preserve every global minimizer in
$B_0$. Their intersection $B_\infty$ is a nonempty box. This elementary
convergence statement does not supply a rate or a computable distance to
$B_\infty$. Nor does nestedness alone justify identifying the intersection
as a fixed point without an additional continuity property of the
relaxation construction. The certificates below avoid that inference.

Write $w(B)=\max_i(u_i-\ell_i)$. A frequently useful bound from the
earlier study illustrates the role of the incumbent and of curvature.

\begin{proposition}[Quadratic growth and a relaxation error]\label{prop:growth}
Suppose a global minimizer $x^*$ and constants $\gamma>0$, $\tau\ge0$
satisfy, for all relevant $B$ containing $x^*$ and all $x\in B$,
\[
 \phi_B(x)\ge f^*+\frac{\gamma}{2}\norm{x-x^*}_\infty^2-\tau w(B)^2.
\]
With $\epsilon=U-f^*\ge0$,
\begin{equation}\label{eq:growth-recurrence}
 w(T_U(B))^2\le\frac{8\epsilon}{\gamma}
                     +\frac{8\tau}{\gamma}w(B)^2.
\end{equation}
If $q^2=8\tau/\gamma<1$, then
\[
 w(B_k)^2\le q^{2k}w(B_0)^2+
 \frac{8\epsilon}{\gamma}\frac{1-q^{2k}}{1-q^2}.
\]
\end{proposition}
\begin{proof}
Every point of $K_U(B)$ is within infinity-norm distance
$\sqrt{2(\epsilon+\tau w(B)^2)/\gamma}$ of $x^*$.
Its coordinate ranges are at most twice this radius. Squaring gives
\eqref{eq:growth-recurrence}; iteration of the affine scalar recurrence
gives the second bound.
\end{proof}

This is a sufficient contraction estimate, not a claim of an exact rate.
It also explains why improving the incumbent may justify reconsidering
tightening: the displayed floor scales with $\sqrt{U-f^*}$. In ordinary
search, $f^*$, $x^*$, and a certified growth constant are not generally
available. The estimate is therefore not itself a computable universal
activation rule. Section~\ref{sec:constraints} treats constrained support
problems and makes the additional sensitivity assumptions explicit.

Finally, relaxation witnesses need not be feasible for the original
nonlinear problem. A witness can prove that a relaxation cannot tighten a
coordinate even when it violates a nonlinear equality or integrality.
Conversely, a feasible incumbent justifies the cutoff but says little
about the remaining width of a weak relaxation. Both types of information
are useful, and they have different proof obligations.
